\title{DeepSeekMath: Pushing the Limits of Mathematical Reasoning in Open Language Models}
\begin{itemize}[topsep=0pt]
    \item The possible influence of ArXiv data on mathematical tasks outside the present evaluation, e.g., the translation of formal mathematical statements or proofs into informal language;
    \item How ArXiv data might affect model performance when combined with other training data sources;
    \item The potential for positive effects from ArXiv papers to appear at even larger model scales.
\end{itemize}

\paragraph{Results}
The empirical outcomes under these varied regimes are presented in Table \ref{tab:code-math} and Table \ref{tab:code-math-general-eval}.

As illustrated in Figure \ref{fig:rl-analysis}, GRPO demonstrates superior performance over online RFT, underscoring the effectiveness of adjusting both positive and negative gradient coefficients. Moreover, GRPO+PS outperforms GRPO+OS, suggesting that employing step-specific and fine-grained gradient coefficients confers additional benefits. Additionally, we investigate the impact of iterative RL by performing two iterations in our experiments. The results presented in Figure \ref{fig:iter-rl} reveal that iterative RL considerably boosts performance, particularly during the first iteration.

\subsubsection{Training Setting}
\label{sec:quality-policy}
We conduct mathematical training using a general-purpose pre-trained language model with 1.3B parameters, utilizing the architecture underlying the DeepSeek LLMs \citep{deepseek-llm}; this configuration is referenced as DeepSeek-LLM 1.3B. For each mathematical corpus, a separate model is trained for a total of 150B tokens. All training is executed via the efficient HAI-LLM \citep{haillm} framework. Adhering to DeepSeek LLMs’ established procedures, we employ the AdamW optimizer \citep{adamW} with hyperparameters $\beta_1=0.9$, $\beta_2=0.95$, and a $\mathrm{weight\_decay}=0.1$. The learning rate follows a multi-stage schedule: it peaks after 2,000 warmup steps, then decreases to 31.6\% of the peak after 80\% of training, and further declines to 10.0\% by 90\% of training completion. The learning rate’s maximum is set at 5.3e-4, the batch size at 4M tokens, and the context length at 4K.

To prevent contamination of evaluation benchmarks, we adopt the method in \cite{deepseek-coder}, eliminating web pages that contain either questions or answers from established English mathematical benchmarks (including GSM8K \citep{gsm8k} and MATH \citep{MATH}) as well as Chinese benchmarks (such as CMATH \citep{wei2023cmath} and AGIEval \citep{agieval}). Our filtering approach removes any segment of text containing a 10-gram string that exactly matches a sub-string in the benchmarks from the training set. For benchmark entries comprising fewer than 10 grams, but at least 3 grams, exact string matching is likewise used to filter out compromised web pages.

\subsection{Validating the Quality of the DeepSeekMath~Corpus}

\begin{itemize}[topsep=0pt]
    \item \textbf{MathPile} \citep{mathpile}: An 8.9B-token collection largely comprising scientific ArXiv papers (over 85\%), curated via a suite of cleaning and filtering rules;
    \item \textbf{ArXiv-RedPajama} \citep{redpajama}: A 28.0B-token dataset of ArXiv LaTeX files, with supplementary content such as preambles, comments, macros, and references omitted.
\end{itemize}

We conduct pre-training experiments to compare the DeepSeekMath~Corpus against other recently released math-specific training datasets:

We systematically evaluate the mathematical proficiency of DeepSeekMath-Base 7B, particularly its effectiveness at generating comprehensive mathematical solutions autonomously, solving problems with the assistance of tools, and engaging in formal theorem verification. In addition to mathematical evaluation, we examine the broader capabilities of the model in natural language understanding, logical reasoning, and programming.

In this formulation, $\varepsilon$ and $\beta$ are hyper-parameters, and the advantage $\hat{A}_{i,t}$ is computed only with respect to the relative rewards of outputs within the sampled group, as will be elaborated upon in subsequent subsections. The use of group-relative advantages in GRPO closely matches the comparative structure of typical reward model training, where preference datasets are composed of output pairs for the same question. Furthermore, instead of integrating the KL penalty within the reward as in PPO, GRPO directly penalizes KL divergence between the current and reference policies in the loss, which simplifies the computation of $\hat{A}_{i,t}$. Distinct from the reward-based KL penalty in (\ref{eq:PPO-reward}), the

Because the value function required by PPO typically matches the policy model in scale, this adds significant resource demands in terms of both memory and computation. Moreover, in RL finetuning, the value function operates as a baseline to reduce variance in advantage estimates. However, within the large language model context, only the reward associated with the final token is usually provided by the reward model, complicating the process of learning an accurate value function at every token. To overcome these difficulties, we introduce Group Relative Policy Optimization (GRPO), as visualized in Figure \ref{fig:grpo}. GRPO eliminates the need for a separate value function by using, as a baseline, the mean reward over a set of sampled outputs responding to the same input question. Concretely, for each question $q$, a batch of outputs $\{o_1, o_2, \cdots, o_G\}$ is sampled from $\pi_{\theta_{old}}$, and the policy model is trained to maximize the following objective:

Upon completing pre-training, we further refine DeepSeekMath-Base through mathematical instruction tuning, utilizing chain-of-thought \citep{cot}, program-of-thought \citep{pot,pal}, and tool-integrated reasoning \citep{tora} datasets. The resulting DeepSeekMath-Instruct 7B model outperforms all other 7B models in its class and reaches competitive performance relative to 70B-scale open-source instruction-tuned models.

We further assess performance in general natural language understanding with MMLU \citep{mmlu}, reasoning via BBH \citep{bbh}, and coding proficiency on HumanEval \citep{codex} and MBPP \citep{mbpp}. According to Table \ref{tab:understanding_reasoning_code}, DeepSeekMath-Base 7B demonstrates marked improvement on MMLU and BBH relative to its predecessor, DeepSeek-Coder-Base-v1.5 \citep{deepseek-coder}, highlighting the advantageous effects of mathematical training for language understanding and reasoning. Simultaneously, continued inclusion of code tokens in ongoing training enables DeepSeekMath-Base 7B to maintain the coding benchmark results of DeepSeek-Coder-Base-v1.5. Taken together, DeepSeekMath-Base 7B substantially exceeds the performance of the generic Mistral 7B \citep{mistral} across all three reasoning and coding evaluation tasks.

    \item In contrast to prevailing practices, our investigation finds that incorporating arXiv papers—commonly used in math-centric corpora—offers no observable benefits across any of the mathematical evaluation benchmarks reported in this study.
\end{itemize}

Furthermore, code pretraining also offers gains in mathematical reasoning tasks that do not involve external computational tools. Within the two-stage scheme, commencing with code training leads to notable, though more modest, gains and improves the efficacy of subsequent math specialization—culminating in the strongest outcomes. By contrast, adopting a one-stage joint training strategy appears to hinder tool-free mathematical reasoning. A plausible explanation is that the relatively compact DeepSeek-LLM 1.3B may lack sufficient model capacity to simultaneously internalize both extensive code and mathematical domains when trained together.

These models are designed for broad applicability, with the majority undergoing various alignment stages.

\subsubsection{Evaluation Results}

\subsubsection{How to Achieve More Effective RL?}
\label{sec:effec_rl}
Our study demonstrates that reinforcement learning exhibits strong effectiveness in tasks involving mathematical reasoning. Furthermore, we propose a unified conceptual framework encompassing various well-known training approaches, framing all such methods as forms or simplifications of RL methodologies. As encapsulated by Equation \ref{eq:objective}, this framework features three fundamental elements: Data Source, Algorithm, and Reward Function. We outline several prospective directions for each of these elements.

\begin{itemize}[topsep=0pt]
    \item We present Group Relative Policy Optimization (GRPO), a novel reinforcement learning algorithm that is both resource-efficient and highly effective. GRPO eliminates the need for a critic network by estimating baselines directly from group-level scores, thereby offering substantial reductions in computational requirements when compared to Proximal Policy Optimization (PPO).
    \item Our experiments show that integrating GRPO into instruction tuning considerably improves the capabilities of our DeepSeekMath-Instruct model, using only the data designated for instruction tuning. In addition, we find that the reinforcement learning phase confers generalization benefits, enhancing out-of-domain performance.
    \item We develop a cohesive framework that unifies approaches such as RFT, DPO, PPO, and GRPO under a common perspective. To thoroughly analyze this paradigm, we undertake a broad array of empirical studies including comparisons between online and offline learning, outcome versus process-based supervision, and single-step versus iterative reinforcement learning, among others.
    \item Leveraging our unified framework, we investigate the underlying mechanisms contributing to the effectiveness of reinforcement learning, and highlight multiple promising directions for advancing the reinforcement learning of large language models (LLMs).
\end{itemize}

\begin{algorithm}[t]
  \small
  \caption{Iterative Group Relative Policy Optimization}
  \textbf{Input} initial policy model $\pi_{\theta_{\text{init}}}$; reward models $r_\varphi$; task prompts $\mathcal{D}$; 
  hyperparameters $\varepsilon$, $\beta$, $\mu$
  \begin{algorithmic}[1]
    \State Set policy model $\pi_\theta \leftarrow \pi_{\theta_{\text{init}}}$
    \For{iteration = 1, \dots, I}
       \State Assign current policy as reference model: $\pi_{ref} \leftarrow \pi_{\theta}$
      \For{step = 1, \dots, M}
      \State Draw a mini-batch $\mathcal{D}_b$ from $\mathcal{D}$
      \State Update prior policy $\pi_{\theta_{old}} \leftarrow \pi_{\theta}$ 
      \State For every question $q \in \mathcal{D}_b$, sample $G$ outputs $\{o_i\}_{i=1}^G \sim \pi_{\theta_{old}} (\cdot \mid q)$
      \State Evaluate each output $o_i$ with the reward model $r_{\varphi}$ to obtain rewards $\{r_i\}_{i=1}^{G}$
      \State Calculate $\hat{A}_{i,t}$, the relative advantage for each token $t$ in $o_i$, using group-based estimation.
      \For{GRPO iteration = 1, \dots, $\mu$}
        \State Adjust $\pi_{\theta}$ to increase the GRPO objective (Equation \ref{eq:GC-GRPO})
      \EndFor
    \EndFor \State Continuously train $r_\varphi$ using a replay buffer.
    \EndFor 
  \end{algorithmic}
  \textbf{Output} $\pi_\theta$
  \label{alg:iter-grpo}
\end{algorithm}

\begin{itemize}[topsep=0pt]
    \item We demonstrate that Common Crawl, a publicly available web corpus, encompasses a rich reservoir of mathematical data. Leveraging a carefully crafted filtering and selection framework, we assemble the DeepSeekMath Corpus—a high-caliber dataset consisting of 120B tokens sourced from web pages specifically curated for mathematical relevance. This corpus exceeds the size of the math web collection utilized by Minerva \citep{minerva} by almost a factor of 7, and surpasses the recently introduced OpenWebMath \citep{openwebmath} by a factor of 9.

While ArXiv papers frequently form a substantial portion of math-focused pre-training datasets \citep{minerva,gpt-f,llemma,mathpile}, their specific influence on the enhancement of mathematical reasoning capabilities has not been rigorously investigated. Somewhat unexpectedly, our empirical findings suggest that the inclusion of ArXiv papers does not lead to measurable improvements in mathematical reasoning. To investigate this, we assessed models of different scales—namely, DeepSeek-LLM 1.3B and DeepSeek-Coder-Base-v1.5 7B \citep{deepseek-coder}—using ArXiv-based datasets processed through varying pipelines:

To evaluate the impact of code pre-training on mathematical reasoning performance, we conducted comparative studies utilizing both two-stage and single-stage training methodologies:

\paragraph{Reward Function} The reward function provides the foundational signal for model training. Within RL, this function is often realized as a neural reward model. We identify three critical avenues for advancing reward models: 1) \textbf{Enhancing the generalization capacity of reward models.} For reinforcement learning to yield meaningful improvements to LLMs, reward models need to generalize robustly to novel or out-of-distribution prompts and sophisticated model outputs—otherwise, RL risks merely reinforcing the current model distribution; 2) \textbf{Quantifying and representing the uncertainty in reward models.} Uncertainty estimates may serve as a crucial interface between inherently weak reward models and weak-to-strong alignment techniques; 3) \textbf{Constructing high-quality, process-level reward models efficiently} that can deliver granular supervisory signals for complex reasoning tasks \citep{lightman2023let,wang2023math}.

    DeepSeekMath-Base achieves performance on par with the proprietary Minerva 540B \citep{minerva} model on English benchmarks, outperforming all open-source baselines—including Mistral 7B \citep{mistral} and Llemma-34B \citep{llemma}—across models regardless of prior mathematical pre-training, often by wide margins. For Chinese benchmarks, DeepSeekMath-Base excels, which we attribute to our inclusion of diverse, high-quality math pre-training data from multiple languages, diverging from the English-centric data curation seen in earlier work \citep{minerva,llemma}. Furthermore, after applying mathematical instruction tuning and reinforcement learning, the models DeepSeekMath-Instruct and DeepSeekMath-RL set new benchmarks for open-source systems, exceeding 50\% accuracy on the competition-level MATH dataset—a first in the open-source domain.
\end{itemize}

\paragraph{Mathematical Problem Solving with Tool Use} 
We assess program-augmented mathematical reasoning abilities using the GSM8K and MATH benchmarks via few-shot program-of-thought prompting \citep{pot,pal}. For each task, models are prompted to produce a Python script that leverages libraries such as \textit{math} and \textit{sympy} for complex operations, with the correctness judged by the program's output. Table \ref{tab:base_math_pot_proof} shows that DeepSeekMath-Base 7B establishes new state-of-the-art results, outperforming the former leading model, Llemma 34B.

Initially, we select OpenWebMath \citep{openwebmath}—a high-quality compilation of mathematics-related online materials—as our foundational seed corpus. Leveraging this dataset, we train a fastText model \citep{joulin2016fasttext} to identify and retrieve additional web pages resembling those in OpenWebMath. Concretely, we draw 500,000 samples from the seed data as positive examples and pair them with 500,000 randomly chosen web pages from Common Crawl as negatives. For model training, we utilize a public fastText library\footnote{\url{https://fasttext.cc}}, specifying a vector size of 256, a learning rate of 0.1, a maximum n-gram length of 3, a minimum word count threshold of 3, and train over 3 epochs. To pare down the volume of Common Crawl, we implement both URL-based deduplication and near-duplicate removal, thereby narrowing the dataset to 40 billion HTML documents. Subsequently, the trained fastText model is used to retrieve mathematical web pages from this deduplicated pool. To ensure high content quality, we score and rank all retrieved pages with the model, retaining only those with the highest scores. The amount of retained data is determined based on pre-training performance across candidate token volumes of 40B, 80B, 120B, and 160B. For the initial round, we preserve the top 40B tokens.

In the evaluation condition permitting both natural language reasoning and code-based tool integration for problem solving, DeepSeekMath-Instruct 7B achieves nearly 60\% accuracy on MATH, outperforming all previous open-source contenders. On other test suites, the model exhibits performance competitive with DeepSeek-LLM-Chat 67B, a previous leading open-source model that is an order of magnitude larger.

The constituent elements of these methods are outlined in Table \ref{tab:data_policy}. For a comprehensive derivation, see Appendix \ref{app:analysis-rl}.

Nevertheless, these findings are subject to important caveats. The following aspects were not addressed in this study:

\paragraph{Insight on Data Source} We categorize data sources into two main types: online sampling and offline sampling. Online sampling refers to acquiring training data from real-time exploration conducted by the actively updated training policy model. Offline sampling, in contrast, sources its training data from the initial SFT model's samples. In this framework, both RFT and DPO are characterized by offline sampling, whereas Online RFT and GRPO employ online sampling approaches.

\textbf{The DeepSeekMath~Corpus demonstrates superior quality, substantial multilingual coverage in mathematics, and unmatched scale.}

Table \ref{tab:sft_rl_math} presents the results achieved by both open- and closed-source models employing chain-of-thought as well as tool-based reasoning approaches on English and Chinese assessment tasks. The key observations are as follows: 1) DeepSeekMath-RL 7B records accuracy rates of 88.2\% on GSM8K and 51.7\% on MATH when using chain-of-thought reasoning. These figures outperform all other open-source models ranging from 7B to 70B parameters, in addition to exceeding the accuracy of most proprietary models. 2) Importantly, DeepSeekMath-RL 7B is exclusively trained on instruction tuning data formatted with chain-of-thought reasoning from GSM8K and MATH, initialized from DeepSeekMath-Instruct 7B. Although its training set is narrowly defined, it consistently surpasses DeepSeekMath-Instruct 7B on every evaluation criterion, indicating the advantage conferred by reinforcement learning.

\subsubsection{Outcome Supervision RL with GRPO}
More precisely, for each input question $q$, a set of $G$ outputs $\{o_1, o_2, \cdots, o_G\}$ are drawn from the preceding policy $\pi_{\theta_{old}}$. Each output is scored by the reward model, producing a reward vector $\mathbf{r}=\{r_1, r_2, \cdots, r_G\}$. These scores are standardized by subtracting the mean of the group and dividing by their standard deviation. Under outcome supervision, the normalized reward is applied uniformly to the final token in each output $o_i$, assigning the advantage for all tokens in $o_i$ as this normalized value: $\hat{A}_{i, t} = \widetilde{r}_i = \frac{r_i- {\rm mean}(\mathbf{r})}{{\rm std}(\mathbf{r})}$. The policy is then improved by optimizing the objective described in equation (\ref{eq:GRPO-obj}).

As depicted in Figure \ref{fig:rl-analysis}, it is observed that Online RFT achieves superior performance compared to RFT across two distinct benchmarks. Initially, the performance of Online RFT aligns closely with that of RFT, but as training progresses, Online RFT secures a clear performance lead, illustrating the benefits of online data collection. This trend is understandable: early in training, the policy (actor) closely mirrors the SFT model, leading to sampled data that diverges only minimally. At later stages, however, the samples generated by the evolving actor diverge more noticeably from the initial model, thus making real-time sampling markedly more advantageous.

\paragraph{Insight on Gradient Coefficient} The algorithm modifies the gradient coefficient based on the processed reward signal to update model parameters. In our study, the reward function is divided into two types: `Rule' and `Model'. The `Rule' method evaluates response quality based on the correctness of answers, while the `Model' method uses a learned reward model to score responses, where this model is trained using data derived from rule-based judgments. Equations \ref{eq:GC-RFT} and \ref{eq:GC-GRPO} demonstrate a salient distinction between GRPO and Online RFT: GRPO dynamically scales its gradient coefficient with the reward assigned by the reward model, enabling nuanced reinforcement or penalization based on the reward's magnitude. Conversely, Online RFT applies uniform reinforcement to all correct answers and does not penalize incorrect ones, as this approach does not incorporate reward-value-based adjustment.

\textbf{Exploration and Analysis of Reinforcement Learning}

Following the initial round of data collection, a significant number of mathematical web pages remain unacquired. This is primarily attributed to the lack of diversity among the positive samples used to train the fastText model. To address this limitation, we expand our pool of seed data by sourcing additional mathematical web domains to better inform the model. We begin by segmenting the entire Common Crawl dataset into mutually exclusive domains, where each domain encompasses web pages sharing the same root URL. For every domain, we determine the proportion of pages captured during the first round of collection. If more than 10\% of a domain's pages are acquired, we categorize it as math-oriented (for instance, \url{mathoverflow.net}).

\subsection{Summary of Evaluations and Metrics}

\subsection{Reflections on Reinforcement Learning}
\subsubsection{Advancing a Unified Analytical Framework} This section proposes an integrated perspective for examining various training approaches—including SFT, RFT, DPO, PPO, GRPO—and presents experimental results that dissect the components of this framework. Broadly, the gradient of a training method with respect to its parameters $\theta$ can be formalized as follows:

As indicated in Table \ref{tab:base_math_cot}, DeepSeekMath-Base 7B consistently achieves top results across all eight evaluated benchmarks among open-source base models—including established models such as Mistral 7B \citep{mistral} and the more recent Llemma 34B \citep{llemma}, which was specifically trained on Proof-Pile-2 \citep{llemma} for mathematical capability. In particular, on the highly challenging MATH dataset, DeepSeekMath-Base attains an improvement exceeding 10\% absolute over all previously available open-source base models, and notably surpasses Minerva 540B \citep{minerva}—a proprietary base model based on PaLM \citep{palm}, over 77 times larger, and additionally trained on specialized mathematical corpora.

We begin by detailing our experiences and observations during pre-training. Unless noted otherwise, the training protocols referenced align with those described in Section \ref{sec:quality-policy}. For clarity, references to the DeepSeekMath Corpus within this discussion pertain to the 89B-token dataset derived from the second round of data collection.

\paragraph{Mathematical Problem Solving with Step-by-Step Reasoning}

\subsection{Data Collection and Decontamination} This section describes the methodology for assembling the DeepSeekMath Corpus using data sourced from Common Crawl. As illustrated in Figure \ref{fig:data_collect}, we introduce an iterative pipeline designed to systematically extract a substantial mathematical dataset from Common Crawl, beginning with a seed collection (e.g., a smaller but rigorously curated set of math datasets). Notably, the outlined strategy is extensible to other specialized areas such as programming.

Here, we present DeepSeekMath-Base 7B, a foundational model notable for its advanced reasoning capacities in mathematics. This model builds upon DeepSeek-Coder-Base-v1.5 7B \citep{deepseek-coder} as its starting point and undergoes training on 500B tokens. The data composition comprises 56\% from the DeepSeekMath Corpus, 4\% from AlgebraicStack, 10\% from arXiv, 20\% drawn from Github code, and the final 10\% derived from Common Crawl natural language data, covering both English and Chinese texts. The training procedure predominantly follows the specifications given in Section \ref{sec:quality-policy}, with modifications including setting the peak learning rate at 4.2e-4 and utilizing a batch size corresponding to 10M tokens.

Our protocol involved individually training DeepSeek-LLM 1.3B (for 150B tokens) and DeepSeek-Coder-Base-v1.5 7B (for 40B tokens) exclusively on each ArXiv dataset. The experimental results indicate that models exposed solely to ArXiv-derived corpora do not achieve noticeable gains—and may in some cases perform worse—on an array of mathematical evaluation tasks varying in difficulty. These include quantitative reasoning assessments like GSM8K and MATH (see Table \ref{tab:arxiv-cot}), STEM-related multiple-choice benchmarks such as MMLU-STEM (Table \ref{tab:arxiv-cot}), and formal mathematical reasoning, for example, miniF2F (refer to Table \ref{tab:arxiv_atp}).

\subsection{Contributions}
This work presents advancements in large-scale math-focused pre-training and delves into the assessment and utilization of reinforcement learning strategies.

\paragraph{Data Source} The data source constitutes the foundational input for any training approach. Within the RL context, this specifically refers to unlabeled questions along with their corresponding model-generated responses. In our experiments, we restrict our data source to questions derived from the instruction tuning phase, with outputs sampled via basic nucleus sampling. We hypothesize that this constraint may be a contributing factor for the observed improvement in Maj@K alone. Moving forward, we plan to extend our RL pipeline to address out-of-distribution prompts and to leverage \textbf{more advanced sampling (decoding) strategies}, including those employing tree-search techniques \citep{yao2023tree}. In addition, \textbf{efficient inference methodologies} \citep{xia-etal-2023-speculative,leviathan2023fast,kwon2023efficient,xia2024unlocking}—which influence the exploration effectiveness of policy models—are likely to be critical to further progress.

\textbf{Math Pre-Training at Scale}

\item \textbf{Formal Mathematics}: We assess DeepSeekMath-Base’s capabilities through the informal-to-formal theorem proving task as described by \citep{dsp_proof}, utilizing the miniF2F benchmark \citep{minif2f} and selecting Isabelle \citep{isabelle} as the proof assistant. In these settings, DeepSeekMath-Base exhibits robust few-shot autoformalization ability.

\paragraph{Algorithms} Algorithms utilize both data and reward signals to calculate the gradient coefficient, thereby facilitating updates to model parameters. According to Equation \ref{eq:objective}, most contemporary approaches essentially \textbf{TRUST} the reward function's signal to modulate the conditional probability of generating specific tokens. Nonetheless, the reliability of the reward signal cannot always be guaranteed, particularly when tackling highly intricate tasks. As an illustration, even rigorously curated datasets like PRM800K \citep{lightman2023let}, which benefit from expert annotator oversight, are estimated to contain roughly 20\% erroneous labels\footnote{\url{https://github.com/openai/prm800k/issues/12\#issuecomment-1728491852}}. Consequently, we propose to investigate reinforcement learning algorithms designed to withstand the presence of noisy reward feedback. In our view, \textbf{WEAK-TO-STRONG} alignment paradigms \citep{burns2023weak} have the potential to revolutionize learning algorithm frameworks at a fundamental level.

The advent of large language models (LLMs) has transformed the landscape of mathematical reasoning within artificial intelligence, leading to marked progress in quantitative reasoning \citep{MATH} as well as in benchmarks focused on geometry \citep{trinh2024solving}. These systems have also demonstrated significant value in aiding human users to address intricate mathematical challenges \citep{tao}. Despite these advances, state-of-the-art solutions such as GPT-4 \citep{gpt4} and Gemini-Ultra \citep{gemini} are not openly released, and the current generation of open-source alternatives lags considerably in terms of performance.

\begin{equation}
\footnotesize
                \mathcal{J}_{PPO}(\theta) = \mathbb{E}{[q \sim P(Q), o \sim \pi_{\theta_{old}}(O|q)]} \frac{1}{|o|} \sum_{t=1}^{|o|} \min \left[ \frac{\pi_\theta(o_{t} | q, o_{<t})}{\pi_{\theta_{old}}(o_{t} | q, o_{<t})} A_{t}, \text{clip} \left( \frac{\pi_\theta(o_{t} | q, o_{<t})}{\pi_{\theta_{old}}(o_{t} | q, o_{<t})}, 1 - \varepsilon, 1 + \varepsilon \right)  A_{t} \right] ,
\end{equation}

\end{document}

\item \textbf{Natural Language Understanding, Reasoning, and Code}: In order to provide a thorough evaluation of the model’s general abilities in comprehension, reasoning, and coding, we benchmark DeepSeekMath-Base on the Massive Multitask Language Understanding (MMLU) dataset \citep{mmlu}, comprising 57 multiple-choice tasks across diverse subject areas; the BIG-Bench Hard (BBH) collection \citep{bbh}, which features 23 highly challenging, multi-step reasoning problems; as well as HumanEval \citep{codex} and MBPP \citep{mbpp}, both widely recognized for assessing code generation models. The results indicate that math-focused pre-training enhances both linguistic understanding and reasoning skills.

\subsubsection{From PPO to GRPO} Proximal Policy Optimization (PPO) \citep{schulman2017proximal} is a widely-adopted actor-critic algorithm for RL-based fine-tuning of LLMs \citep{ouyang2022training}. PPO seeks to optimize LLMs by maximizing the following objective function:

\subsubsection{Code Training Enhances Mathematical Reasoning}

DeepSeekMath-Base is built upon DeepSeek-Coder-Base-v1.5 7B \citep{deepseek-coder}, as we determined that leveraging a model initially trained on code offers advantages over beginning with a standard LLM. Notably, we observe that mathematical training not only sharpens the model’s proficiency in mathematics but also leads to performance gains on benchmarks such as MMLU \citep{mmlu} and BBH \citep{bbh}, suggesting an overall improvement in broader reasoning skills.

\subsection{Group Relative Policy Optimization} Reinforcement learning (RL) has been demonstrated to further bolster the mathematical reasoning capabilities of LLMs following Supervised Fine-Tuning (SFT) \citep{wang2023math,wizardmath}. In this part, we present Group Relative Policy Optimization (GRPO), an efficient and robust RL strategy.

\section{Conclusion, Limitation, and Future Work} This work introduces DeepSeekMath, a model that surpasses all other open-source systems on the MATH benchmark at the competition level, and closely approaches the results of proprietary models. DeepSeekMath~begins from the DeepSeek-Coder-v1.5 7B checkpoint and is subjected to continuous training for 500B tokens, including a substantial portion—120B tokens—of mathematical content extracted from Common Crawl. Through a comprehensive ablation analysis, we find that web data is a valuable source for quality mathematical material, whereas arXiv contributes less than anticipated. We propose Group Relative Policy Optimization (GRPO), which modifies the standard Proximal Policy Optimization (PPO) framework, and demonstrates significant gains in mathematical reasoning with reduced memory requirements. Experimental evidence suggests that GRPO enhances model ability even when DeepSeekMath-Instruct 7B is already achieving high benchmark scores. Additionally, we outline a cohesive framework for interpreting various reinforcement learning approaches and identify several promising research directions for improving RL efficacy.

We assemble a diverse instruction-tuning dataset for mathematics that includes both English and Chinese problems. The dataset spans a variety of mathematical disciplines and complexity levels. Each problem is paired with corresponding solutions that follow different reasoning formats, namely chain-of-thought (CoT) \citep{cot}, program-of-thought (PoT) \citep{pot,pal}, and tool-integrated reasoning approaches \citep{tora}. In total, this training corpus consists of 776K samples.

\subsubsection{Process Supervision RL with GRPO}
Although outcome supervision only gives a terminal reward for each output, such sparse feedback can be insufficient for guiding policies on tasks requiring multi-step reasoning, such as mathematics. Building on \cite{wang2023math}, we examine process supervision, where feedback is provided at every intermediate reasoning step. Formally, for a given question $q$ and $G$ sampled completions $\{o_1, o_2, \cdots, o_G\}$, a process reward model produces scores for each reasoning stage, yielding $\mathbf{R} = \{ \{r_1^{index(1)},\cdots,r_1^{index(K_1)}\}, \cdots,  \{r_G^{index(1)},\cdots,r_G^{index(K_G)}\} \}$, with $index(j)$ denoting the terminal token of the $j$-th step, and $K_i$ the number of reasoning steps in output $o_i$. These per-step rewards are also normalized: $\widetilde{r}_i^{index(j)} = \frac{r_i^{index(j)} - {\rm mean(\mathbf{R})}}{{\rm std(\mathbf{R})}}$.
For each token $t$, process supervision computes its advantage as the cumulative sum of normalized rewards from subsequent steps, specifically $\hat{A}_{i, t} = \sum_{index(j) \ge t} \widetilde{r}_i^{index(j)}$. The policy is refined by maximizing the same objective as in equation (\ref{eq:GRPO

\section{Reinforcement Learning}

\section{Introduction}

Additionally, we propose Group Relative Policy Optimization (GRPO), an RL variant derived from Proximal Policy Optimization (PPO) \citep{schulman2017proximal}. GRPO eliminates the need for a critic network by estimating the baseline directly from group score distributions, thus reducing computational overhead during training. When applied to a selected subset of English instruction tuning examples, GRPO yields notable gains over the robust DeepSeekMath-Instruct baseline, both in-domain (e.g., GSM8K: 82.9\% $\rightarrow$ 88.2\%, MATH: 46.8\% $\rightarrow$ 51.7\%) and on out-of-domain math tasks (such as CMATH: 84.6\% $\rightarrow$ 88.8\%) during reinforcement learning.

\subsection{Training and Evaluating DeepSeekMath-RL}

where $r_\varphi$ denotes the output of the reward model, $\pi_{ref}$ stands for the reference policy (typically the base SFT model), and $\beta$ scales the strength of the KL penalty term.

\paragraph{Formal Mathematics}

\begin{itemize}[topsep=0pt]
    \item \textbf{MathPile} \citep{mathpile}: a multi-source dataset (8.9B tokens), drawn from textbooks, Wikipedia, ProofWiki, CommonCrawl, StackExchange, and arXiv, where more than 85\% of the content originates from arXiv;
    \item \textbf{OpenWebMath} \citep{openwebmath}: mathematical content extracted from CommonCrawl, totaling 13.6B tokens;
    \item \textbf{Proof-Pile-2} \citep{llemma}: a mathematical compilation including OpenWebMath, AlgebraicStack (with 10.3B tokens of mathematical code), and arXiv literature (28.0B tokens). For experiments using Proof-Pile-2, we apply the arXiv:Web:Code ratio of 2:4:1 as specified in \cite{llemma}.
\end{itemize}

    \item We present experimental insights from our math pre-training regimen. Initial training on code datasets prior to mathematical data exposure results in enhanced mathematical problem-solving proficiency, regardless of whether external tools are available. This observation lends evidence towards answering the enduring query: \textit{does training on code promote better reasoning?} Our results suggest a positive effect, particularly in mathematical reasoning contexts.

\subsection{Training and Evaluating DeepSeekMath-Instruct 7B}

\section{Discussion}
Here, we present insights gained from both our pre-training and reinforcement learning experiments.

As indicated in Table \ref{tab:sft_rl_math}, when evaluation excludes tool use, DeepSeekMath-Instruct 7B excels in performing multi-step reasoning. Impressively, on the MATH dataset—known for its competition-level challenges—our model outperforms all currently available open-source and nearly all proprietary systems (including Inflection-2 and Gemini Pro) by a margin of at least 9\% in absolute terms. This holds true even relative to models that are significantly larger (for example, Qwen 72B) or those that have received math-targeted reinforcement learning enhancements (such as WizardMath-v1.1 7B). Although DeepSeekMath-Instruct performs on par with the Chinese closed models GLM-4 and Baichuan-3 in MATH evaluations, it still lags behind state-of-the-art proprietary models like GPT-4 and Gemini Ultra.

\begin{itemize}
    \item  \textbf{Supervised Fine-tuning (SFT)}: SFT involves updating a pretrained model using curated human-annotated data.
    \item \textbf{Rejection Sampling Fine-tuning (RFT)}: RFT enhances the SFT model by further training on samples generated by the SFT model, retaining only those responses deemed correct according to an answer-checking criterion.
    \item \textbf{Direct Preference Optimization (DPO)}: DPO iteratively improves upon the SFT baseline by fine-tuning on expanded sets of outputs from the SFT model, optimizing a pairwise loss based on preference signals.
    \item \textbf{Online Rejection Sampling Fine-tuning (Online RFT)}: In contrast to standard RFT, Online RFT begins with the SFT model as a starting point but continually updates the policy using outputs sampled from the actively trained policy model, reflecting its current state.
    \item \textbf{PPO/GRPO}: PPO/GRPO leverages the SFT model for initialization and iteratively strengthens the policy by reinforcing it through real-time sampled outputs.
\end{itemize}

\begin{equation}
\footnotesize
\begin{split}
    \mathcal{J}_{GRPO}(\theta) &= \mathbb{E}{[q \sim P(Q), \{o_i\}_{i=1}^G \sim \pi_{\theta_{old}}(O|q)]}  \\
    & \frac{1}{G}\sum_{i=1}^G\frac{1}{|o_i|} \sum_{t=1}^{|o_i|} \left\{ \min \left[ \frac{\pi_\theta(o_{i,t} | q, o_{i,<t})}{\pi_{\theta_{old}}(o_{i,t} | q, o_{i,<t})} \hat{A}_{i,t}, \text{clip} \left( \frac{\pi_\theta(o_{i,t} | q, o_{i,<t})}{\pi_{\theta_{old}}(o_{i,t} | q, o_{i,<t})}, 1 - \varepsilon, 1 + \varepsilon \right)  \hat{A}_{i,t} \right] - \beta \mathbb{D}_{KL}\left[\pi_{\theta} || \pi_{ref}\right]\right\} ,
\end{split}
\label{eq:GRPO-obj}
\end{equation}

\begin{equation}
    r_{t} = r_\varphi(q, o_{\le t}) - \beta \log\frac{\pi_{\theta}(o_{t}|q, o_{<t})}{\pi_{ref}(o_{t}|q, o_{<t})},
\label{eq:PPO-reward}
\end{equation}

This formulation comprises three principal elements: 1) \textit{Data Source $\mathcal{D}$}, which specifies the collection of data for training; 2) \textit{Reward Function $\pi_{{rf}}$}, responsible for providing the evaluative feedback during learning; 3) \textit{Algorithm $\mathcal{A}$}, which utilizes both the data and reward information to compute the gradient coefficient $GC$, thereby determining the strength of reward or penalty applied. We discuss several canonical algorithms within this general framework:

\textbf{One-Stage Training}

Here, $\pi_{\theta}$ represents the current policy model, while $\pi_{\theta_{old}}$ denotes the previous policy model. The variables $q$ and $o$ correspond to questions and outputs, respectively, with samples drawn from the question dataset and generated by the old policy $\pi_{\theta_{old}}$. The hyper-parameter $\varepsilon$ controls clipping, which is incorporated in PPO to ensure stable training dynamics. The advantage term, $A_t$, is estimated through Generalized Advantage Estimation (GAE) \citep{gae}, utilizing rewards $\{r_{\ge t}\}$ and a parameterized value function $V_{\psi}$. Consequently, PPO requires simultaneous training of a value function with the policy model. To prevent the reward model from being excessively exploited, a common practice involves augmenting the reward at each token with a KL penalty, comparing the trained model to a fixed reference (typically the initial SFT model), as proposed by \citet{ouyang2022training}. The reward function is thus defined as:

Automated formal proof generation contributes to increased rigor, accuracy, and efficiency in mathematical proof processes, and has drawn considerable research interest. We examine DeepSeekMath-Base 7B’s ability in the informal-to-formal proof translation task introduced by \citep{dsp_proof}, where a model generates a formal proof from a given informal theorem statement, its formalized version, and an informal proof sketch. For evaluation, we employ the miniF2F benchmark \citep{minif2f}, which consists of Olympiad-level formal mathematics problems, prompting the model with few-shot examples to produce formal proofs in Isabelle. Adopting the methodology of \cite{dsp_proof}, proof sketches from the model are further processed using the automated prover Sledgehammer \citep{sledgehammer} to complete the proof details. As presented in Table \ref{tab:base_math_pot_proof}, DeepSeekMath-Base 7B achieves strong results in autoformalization.

\begin{itemize}[topsep=0pt]
    \item \textbf{High-quality}:
    We measure downstream capabilities across 8 mathematical benchmarks using few-shot chain-of-thought prompting \cite{cot}. Table \ref{tab:corpora_comparison} reveals that models fine-tuned on the DeepSeekMath~Corpus outperform those trained on alternative corpora. Furthermore, Figure \ref{fig:corpora_comparisons} illustrates the superior results obtained from the DeepSeekMath Corpus compared to Proof-Pile-2 after processing 50B tokens (equivalent to one epoch of Proof-Pile-2), indicating a higher average quality in the DeepSeekMath Corpus. 
    \item \textbf{Multilingual}:
    DeepSeekMath~Corpus incorporates multilingual mathematical content, with English and Chinese constituting its two primary languages. As evidenced in Table \ref{tab:corpora_comparison}, models trained on DeepSeekMath~Corpus exhibit improved mathematical reasoning performance in both English and Chinese contexts. In comparison, conventional mathematical corpora, which predominantly focus on English, display only marginal or even detrimental effects on Chinese mathematical reasoning abilities.
    \item \textbf{Large-scale}:
    The scale of DeepSeekMath~Corpus far exceeds that of prior mathematical corpora. As illustrated in Figure \ref{fig:corpora_comparisons}, training DeepSeek-LLM 1.3B with DeepSeekMath~Corpus results in a more accelerated and sustained learning curve. Conversely, the baseline corpora are much smaller, have undergone numerous training cycles, and lead to models whose performance plateaus significantly earlier.
\end{itemize}

\subsubsection{Limited Effectiveness of ArXiv Papers for Enhancing Mathematical Reasoning}

Despite its strengths on quantitative benchmarks, DeepSeekMath~lags behind closed-source counterparts on tasks involving geometry and theorem-proving. For example, the model is unable to solve problems concerning triangles or ellipses during preliminary testing, which hints at possible bias in data selection during the pre-training and fine-tuning stages. Model size further constrains performance, as DeepSeekMath~does not match GPT-4 in few-shot learning—whereas GPT-4 exhibits notable improvements with few-shot prompts, DeepSeekMath~performs similarly regardless of prompt setting. Our future work will focus on refining the data selection pipeline to curate a higher-quality training corpus, and further investigate the strategies outlined in Section \ref{sec:effec_rl} for advancing reinforcement learning techniques for LLMs.

\subsubsection{Iterative RL with GRPO} During the reinforcement learning procedure, earlier versions of the reward model may become inadequate for effectively supervising the evolving policy model. To address this, we investigate an iterative RL approach utilizing GRPO. As detailed in Algorithm \ref{alg:iter-grpo}, iterative GRPO involves constructing updated training datasets for the reward model using samples generated by the current policy model. The reward model is incrementally improved through training with a replay buffer that includes 10\% of prior data. Subsequently, the updated policy model serves as the new reference model, which is then further optimized using feedback from the updated reward model.

\section{Supervised Fine-Tuning}

\subsection{SFT Data Curation}

Model performance in mathematical reasoning is assessed both in the absence and presence of tool support, across four quantitative reasoning benchmarks in both English and Chinese. For comparison, we include evaluations against several state-of-the-art models available at the time:

We also establish a unified perspective encompassing various approaches, including Rejection Sampling Fine-Tuning (RFT) \citep{yuan2023scaling}, Direct Preference Optimization (DPO) \citep{dpo}, PPO, and GRPO. Within this unified framework, we show that these methods can be interpreted as different forms of direct or streamlined RL algorithms. Our comprehensive experiments—covering comparisons such as online versus offline training, outcome-based versus process supervision, and single-turn versus iterative RL—allow us to dissect critical components of this paradigm. Finally, we clarify the mechanisms through which our RL framework elevates instruction-tuned model performance and outline promising directions for further enhancing RL efficacy within this unified schema.

This paper presents DeepSeekMath, a language model tailored specifically to mathematical tasks that surpasses the capabilities of existing open-source models and approaches the proficiency of GPT-4 on established academic evaluations. The foundation of this model is the DeepSeekMath Corpus, an extensive and meticulously curated pre-training dataset containing 120B mathematics-related tokens. The corpus is compiled from the Common Crawl (CC) dataset utilizing a classifier based on fastText \citep{joulin2016fasttext}. Initially, the classifier is developed with OpenWebMath examples \citep{openwebmath} as positive samples and a diverse range of unrelated web pages as negative samples. The trained classifier is then employed to identify additional relevant entries within the CC, which are subsequently validated and refined through manual annotation. The classifier is retrained using this improved dataset to enhance its accuracy. Our assessment shows that the resulting corpus achieves high standards of quality: the DeepSeekMath-Base 7B model secures scores of 64.2\% on GSM8K \citep{gsm8k} and 36.2\% on the MATH dataset \citep{MATH}, outperforming Minerva 540B \citep{minerva}. Furthermore, owing to the multilingual nature of the DeepSeekMath Corpus, we observe notable progress on Chinese mathematical tasks \citep{wei2023cmath,agieval}. We regard our methodology for mathematical data construction as a preliminary contribution to the field, recognizing ample opportunity for continued advancement.

\item \textbf{Open-source models} encompass: general-purpose systems such as (1) DeepSeek-LLM-Chat 67B \citep{deepseek-llm}, (2) Qwen 72B \citep{qwen}, (3) SeaLLM-v2 7B \citep{seallm}, and (4) ChatGLM3 6B \citep{chatglm3}. Additionally, several models feature advancements tailored for mathematics, including (5) InternLM2-Math 20B~\footnote{\url{https://github.com/InternLM/InternLM-Math}}, a variant of InternLM2 that has been subject to specialized math pretraining and subsequent instruction tuning; (6) Math-Shepherd-Mistral 7B, which uses PPO training \citep{schulman2017proximal} on Mistral 7B \citep{mistral} leveraging a reward model based on supervised processes; (7) the WizardMath series \citep{wizardmath}, which enhances the mathematical reasoning capacity of Mistral 7B and Llama-2 70B \citep{llama2} through evolve-instruct—a form of instruction tuning employing AI-generated prompts—and PPO, with training tasks predominantly drawn from GSM8K and MATH; (8) MetaMath 70B \citep{metamath}, produced by fine-tuning Llama-2 70B on enriched versions of the GSM8K and MATH datasets; (9) ToRA 34B \cite{tora}, which fine-tunes CodeLlama 34B for tool-based mathematical reasoning; and (10) MAmmoTH 70B \citep{MathInstruct}, created by instruction tuning Llama-2 70B on MathInstruct.

\begin{itemize}[topsep=0pt]
    \item \textbf{English mathematical datasets}: For the English data, we provide tool-integrated solutions to problems from GSM8K and MATH, and select samples from MathInstruct \citep{MathInstruct} as well as the Lila-OOD training set \citep{lila}, where problems are solved via CoT or PoT formats. These English datasets encapsulate a broad spectrum of mathematical areas, such as algebra, probability, number theory, calculus, and geometry.
    \item \textbf{Chinese mathematical datasets}: Our Chinese data is composed of K-12 level math problems encompassing 76 distinct sub-topics—including linear equations—each supplemented with solutions presented in CoT and tool-integrated reasoning formats.
\end{itemize}

\textbf{Two-Stage Training}

\begin{itemize}[topsep=0pt]
    \item \textbf{Code Training for 400B Tokens $\rightarrow$ Math Training for 150B Tokens}: The DeepSeek-LLM 1.3B model undergoes pretraining on 400B tokens of code data, followed by an additional 150B tokens of math-specific content;
    \item \textbf{General Training for 400B Tokens $\rightarrow$ Math Training for 150B Tokens}: For comparative purposes, we conduct a control experiment by substituting the initial code data with 400B tokens of general domain text, drawn from DeepSeek-AI’s extensive general corpus. This allows us to systematically explore whether code-centric pretraining offers distinct advantages over generic text in enhancing the model’s mathematical reasoning skills.
\end{itemize}

Introducing code data during pretraining consistently augments the model’s proficiency in program-supported mathematical reasoning tasks, as evidenced in both sequential (two-stage) and mixed (one-stage) training approaches. Notably, Table \ref{tab:code-math} indicates that an initial phase of code training alone markedly strengthens performance on benchmarks such as GSM8K and MATH when solved via Python. Supplementary math-focused training in the second stage yields additional improvements. Moreover, when code and math data are combined for simultaneous one-stage training, the model largely overcomes the catastrophic forgetting observed in strict two-stage protocols, thereby enhancing both coding capability (Table \ref{tab:code-math-general-eval}) and program-mediated math reasoning (Table \ref{tab:code-math}).

There exists a widely held, yet largely unsubstantiated, belief that training models on code can bolster their reasoning skills. Our investigation seeks to shed light on this claim in the context of mathematics: specifically, we observe that exposure to code augments a model's mathematical reasoning capacities, both independently and when leveraging external tools.

\begin{equation}
    \nabla_{\theta}\mathcal{J}_{\textcolor{red}{\mathcal{A}}}(\theta) = \mathbb{E}[\underbrace{(q,o) \sim \textcolor{red}{\mathcal{D}}}_{Data \ Source}]\left( \frac{1}{|o|} \sum_{t=1}^{|o|}  \underbrace{GC_{{\mathcal{A}}}(q, o, t, \textcolor{red}{\pi_{{rf}}})}_{Gradient \ Coefficient}  \nabla_{\theta}\log \pi_{\theta}(o_t | q, o_{<t})\right).
\label{eq:objective}
\end{equation}

\begin{itemize}[topsep=0pt]
    \item \textbf{Math Training for 150B Tokens}: The model is directly trained from scratch using 150B tokens composed solely of mathematical content;
    \item \textbf{Training on a mixture of 400B Code Tokens and 150B Math Tokens}: Since sequential math finetuning on a model pretrained with code often reduces code task performance, we evaluate whether integrating code and math tokens concurrently during a single training phase can bolster mathematical reasoning while also alleviating the detrimental effects of catastrophic forgetting.
\end{itemize}

    \item Our pre-trained model, DeepSeekMath-Base 7B, reaches a performance level on par with the much larger Minerva 540B \citep{minerva}, illustrating that model size is not the sole determinant of mathematical reasoning strength. High-quality data for pre-training can enable smaller models to exhibit competitive mathematical abilities.

\begin{itemize}[topsep=0pt]
    \item \textbf{English and Chinese Mathematical Reasoning}: Our evaluation spans a diverse set of English and Chinese benchmarks that encompass math problems from elementary through collegiate levels. For English, benchmarks include GSM8K \citep{gsm8k}, MATH \citep{MATH}, SAT \citep{llemma}, OCW Courses \citep{minerva}, and MMLU-STEM \citep{mmlu}; for Chinese, we assess MGSM-zh \citep{mgsm}, CMATH \citep{wei2023cmath}, Gaokao-MathCloze \citep{agieval}, and Gaokao-MathQA \citep{agieval}. Model evaluations are based on their abilities to generate comprehensive, independent textual solutions as well as to utilize Python for problem-solving.

\subsubsection{Why RL Works?} In this work, we apply reinforcement learning on a portion of the instruction tuning dataset, resulting in marked improvements relative to the instruction tuning baseline. To further elucidate the efficacy of reinforcement learning, we measure the Pass@K and Maj@K accuracy metrics for both the Instruct and RL-optimized models across two evaluation datasets. As depicted in Figure \ref{fig:maj-pass}, RL increases Maj@K accuracy, yet leaves Pass@K largely unchanged. These observations imply that RL's performance gain primarily arises from enhancing the robustness of the output distribution—specifically, \textbf{the gains appear to derive from an increased likelihood of selecting the correct response from TopK candidates, rather than from improvements to intrinsic model competency.} A similar \textbf{misalignment problem} in reasoning-related SFT models was reported in \citep{wang2023making}, wherein a series of preference alignment techniques \citep{yuan2023rrhf,song2023preference,wang2023making} were shown to bolster the reasoning abilities of SFT models.

\subsection{Training and Evaluating DeepSeekMath-Base 7B}

To assess DeepSeekMath-Base's capacity for mathematical problem solving, we employ few-shot chain-of-thought prompting \citep{cot} across eight benchmark datasets in both English and Chinese. These datasets feature quantitative reasoning tasks (such as GSM8K \citep{gsm8k}, MATH \citep{MATH}, and CMATH \citep{wei2023cmath}) and multiple-choice assessments (including MMLU-STEM \citep{mmlu} and Gaokao-MathQA \citep{agieval}), spanning a wide spectrum of mathematical areas from basic to advanced college-level topics.

\paragraph{Natural Language Understanding, Reasoning, and Code}

\section{Math Pre-Training}

Therefore, more comprehensive investigation is warranted, which we propose as a direction for forthcoming work.

\subsection{Insights from Pre-Training}


\begin{abstract}
The intricacies and systematic characteristics of mathematical reasoning present notable obstacles for language models. This work presents DeepSeekMath~7B, which extends the pre-training of DeepSeek-Coder-Base-v1.5 7B by incorporating 120B mathematics-focused tokens, gathered from Common Crawl alongside data comprising both natural language and code. Remarkably, DeepSeekMath~7B secures a 51.7\% score on the competitive MATH benchmark, managing this without dependence on supplementary toolkits or ensemble voting methods, thus nearing the achievements of models like Gemini-Ultra and GPT-4. When employing self-consistency across 64 outputs, DeepSeekMath~7B attains 60.9\% on the MATH benchmark. The advanced mathematical reasoning exhibited by DeepSeekMath~can be ascribed to two principal factors: Firstly, a rigorously constructed data curation pipeline effectively utilizes substantial volumes of publicly available web data. Secondly, a newly proposed algorithm, Group Relative Policy Optimization (GRPO)—a modified version of Proximal Policy Optimization (PPO)—bolsters mathematical reasoning while also refining PPO's memory efficiency.
\end{abstract}

\begin{itemize}[topsep=0pt]
    \item \textbf{Closed-source models} consist of: (1) members of the GPT model family, with GPT-4 \citep{gpt4} and GPT-4 Code Interpreter~\footnote{\url{https://openai.com/blog/chatgpt-plugins\#\#code-interpreter}} representing the top performers, (2) Gemini Ultra and Pro \citep{gemini}, (3) Inflection-2 \citep{inflection-2}, (4) Grok-1~\footnote{\url{https://x.ai/model-card}}, as well as several Chinese-developed models including (5) Baichuan-3~\footnote{\url{https://www.baichuan-ai.com}}, and (6) the most recent GLM-4~\footnote{\url{https://open.bigmodel.cn/dev/api\#glm-4}} from the GLM family \citep{glm}.
\end{itemize}

Our reinforcement learning experiments are conducted using DeepSeekMath-Instruct 7B as the base model. The RL training set comprises approximately 144K chain-of-thought-formatted questions taken from the GSM8K and MATH subsets within the SFT data, excluding all other SFT questions. This selective dataset construction allows us to analyze the influence of RL in benchmarks without RL exposure. The reward model’s training set is generated following the procedure outlined in \citep{wang2023math}. The initial reward model is fine-tuned from DeepSeekMath-Base 7B using a learning rate of 2e-5. For training with GRPO, we fix the policy model's learning rate at 1e-6, and set the KL divergence penalty coefficient to 0.04. We draw $64$ samples for each prompt, set a maximum sequence length of 1024 tokens, and use a batch size of 1024 during training. The policy model is updated just once after each sampling round. For evaluation, we apply the same benchmarks as those used with DeepSeekMath-Instruct 7B. Under this evaluation scheme, tasks like GSM8K and MATH featuring chain-of-thought reasoning are categorized as in-domain, while remaining benchmarks are treated as out-of-domain.

Next, within these math-related domains, we manually annotate URLs that are known to contain mathematical material (for example, \url{mathoverflow.net/questions}). Uncollected web pages associated with these URLs are then incorporated into the seed corpus. This process allows us to curate a broader range of positive examples and to refine the fastText model, thereby enhancing its retrieval capabilities for mathematical content in subsequent iterations. Through four rounds of this iterative procedure, we assemble a total of 35.5M mathematical web pages, amounting to 120B tokens. By the fourth round, we observe that approximately 98\% of the data had already been retrieved in the preceding round, prompting us to conclude the collection process.

This section details the fine-tuning of DeepSeekMath-Instruct 7B through mathematical instruction optimization, building on DeepSeekMath-Base. Training data examples are concatenated at random until the total token length reaches the model’s context limit of 4K tokens. The model is trained for 500 steps, utilizing a batch size of 256 and maintaining a constant learning rate of 5e-5 throughout.